\subsection*{相同的超时可以来自不同的执行历史}

现在把故障位置放开，只保留客户端的观察：它发出 \code{paper-001}，等候期限耗尽，尚未取得业务结果。下面的每条历史都从余额 10000、业务记录为空开始，各自独立演算。

历史 A 中，请求没有到达服务。连接可能尚未建立，也可能在传输途中中断。服务没有受理 R1，没有生成 P1，数据库也没有开始购买事务。客户端超时以后，账户仍为 10000，订单和流水为空。随后重新发送并成功处理，会形成一笔 3000 分的购买，余额变为 7000。

历史 B 中，服务已经受理请求，在事务中完成了部分写入。它可能已经把余额临时改为 7000，也可能已经写入订单。随后某条语句出错，或数据库最终结束仍然打开的事务，整笔事务回滚。最终可读取的状态恢复为 10000 和空业务历史。

要观察这条回滚路径，可以在订单写入后、流水写入前触发一次 SQL 语句错误，例如除零错误。错误发生前，同一事务内能读到余额 7000 以及购买、订单各一条；回滚完成后，独立连接读到余额 10000，业务记录为空。这里利用的是未提交事务的整体撤销；数据库进程崩溃与断电恢复还涉及各自的恢复机制。

历史 C 是本章开头的情况：事务完整提交，余额为 7000，P1、O1、D1 存在，响应未到客户端。重新 POST 会创建第二笔购买；余额足够时，再扣 3000，得到 4000。

历史 D 中，请求已经到达服务，处理却被延迟。客户端超时那一刻，事务仍未结束，或购买逻辑还在等待执行。原处理可能稍后提交，也可能稍后回滚。客户端这时另发请求，两次处理就可能先后成功。

\begin{figure}[htbp]
\centering
\begin{tikzpicture}[x=1cm,y=1cm]
\node[booknode,text width=10cm] at (5,0) {每条历史的起点：余额 10000，购买、订单、流水为空};
\foreach \i/\y/\a/\b/\c in {
1/-1.4/{A：请求未到服务}/{没有开始购买事务}/{10000；0 笔},
2/-2.8/{B：事务部分执行}/{最终整体回滚}/{10000；0 笔},
3/-4.2/{C：提交后失去响应}/{P1、O1、D1 已保存}/{7000；1 笔},
4/-5.6/{D：仍在途或执行中}/{稍后提交或回滚}/{尚无终局}} {
  \node[booknode,text width=3.2cm] (a\i) at (.1,\y) {\a};
  \node[booknode,text width=3.6cm] (b\i) at (5,\y) {\b};
  \node[booknode,text width=2.5cm] (c\i) at (10,\y) {\c};
  \draw[bookarrow] (a\i)--(b\i);
  \draw[bookarrow] (b\i)--(c\i);
}
\node[booknote,text width=11cm] at (5,-6.6) {客户端都可能等待超时；右列分别统计该分支中的余额与已提交成功购买数};
\end{tikzpicture}
\caption{一次超时覆盖四类不同远端历史。每行独立从相同初始状态推演。}
\label{fig:ch02-four-histories}
\end{figure}


客户端得到的一个超时结果，没有携带这些远端事件。某些错误可以缩小判断范围：例如连接建立前失败，并且能够排除请求曾经传输，就可能支持原请求未被处理的判断。等待响应超时或提交附近断连，覆盖的历史通常更多。

还可以把请求停在 HTTP 入口已经收到报文、尚未调用购买逻辑的位置，保持连接直到客户端超时，再丢弃请求。它与提交后丢弃响应都能使客户端等不到响应头，数据库却分别保留 10000 与空记录、7000 与完整购买。请求在网络中完全未到达是更早的中断位置；它们共同说明，客户端的一种超时观察容得下不同的业务历史。

\subsection*{给不可区分的判断写清前提}

可以把刚才的直觉整理成一个很小的论证。讨论范围限于现有购买入口，以及客户端在有限等待结束时作出的判断。客户端与服务通过消息通信；在这段等待区间内，消息和服务执行可以延迟到截止时间以后，响应也可以丢失。服务可能在执行前或提交后崩溃。客户端尚无另一个可靠通道读取数据库，也没有已确认送达的回滚证明。

这些前提只要求某些消息在有限区间内迟到或丢失。系统可以在后来恢复通信，服务也可以重新可用。

把客户端截至超时观察到的事件序列记为 $L$。它包括发送的请求、本地时钟推进，以及尚未收到完整业务响应这一事实。如果客户端还观察到连接关闭或某项错误，也应计入 $L$。考虑两条远端历史：

\begin{BookCode}
H0：客户端发送 → 请求未进入业务执行 → 客户端超时
H1：客户端发送 → 服务受理并提交 P1 → 响应丢失 → 客户端超时
\end{BookCode}

H0 结束时，成功购买数为 0；H1 结束时，成功购买数为 1。两条历史中，客户端的输入、截止时间和截至这一点的可见事件可以安排成相同。若客户端采用确定性规则，也就是在输入与已经观察到的事件相同时总作出相同选择，那么一个只依据 $L$ 作判断的客户端，就会给出同一个判断。

若它报告未购买，H1 中这个判断错误；若它报告已经购买，H0 中这个判断错误。结果未知与这两种可能性都相容。

\begin{figure}[htbp]
\centering
\begin{tikzpicture}[x=1cm,y=1cm]
\node[bookresult,text width=11.7cm] (view) at (5.2,0)
 {客户端观察 $L$：发送相同请求 $\to$ 等待 $\to$ 到达相同截止时间};
\node[booknode,text width=4.8cm] (h0) at (1.6,-1.7) {H0：请求未进入业务执行\\没有购买提交};
\node[booknode,text width=4.8cm] (h1) at (8.8,-1.7) {H1：购买已经提交\\响应未到客户端};
\draw[bookarrow] (view.south)--++(0,-.35)-|(h0.north);
\draw[bookarrow] (view.south)--++(0,-.35)-|(h1.north);
\node[booknode,text width=4.8cm] (s0) at (1.6,-3.2) {余额 10000\\成功扣款 0 次};
\node[bookresult,text width=4.8cm] (s1) at (8.8,-3.2) {余额 7000\\成功扣款 1 次};
\draw[bookarrow] (h0)--(s0);
\draw[bookalert] (h1)--(s1);
\node[booknote,text width=11.6cm] at (5.2,-4.5) {只以同一份 $L$ 作出同一个二值结论，会在其中一条历史中出错\\查询或后续响应可以增加证据，改变此时的判断依据};
\end{tikzpicture}
\caption{客户端的观察相同，业务效果不同。不可区分只针对当前观察点可用的信息。}
\label{fig:ch02-indistinguishable}
\end{figure}


客户端观察不可区分，指的就是这个观察点上的信息相同。远端历史实际不同；后来到达的响应、数据库查询或恢复记录，仍可能让调用方分清它们。增加了新的证据以后，判断依据也随之改变。

随机选择重试，同样没有增加远端证据。客户端可以按某个概率决定是否重新发送，但在两条历史的局部观察相同时，选择分布也相同。这样的策略可以用于权衡平均损失与等待成本；要让每条允许历史中的确定结论都正确，还需要取得与原操作结果有关的信息。

\subsection*{盲目重试为什么无法同时给出两个保证}

沿用上面的前提，再加上当前接口规则：每次新 POST 可以创建一笔新购买，服务尚无跨请求识别同一操作的机制，余额足够支付两次。客户端希望仅靠超时以后是否重发，同时保证这次购买最终至少生效一次、最多生效一次。

在 H0 中，第一次请求没有产生业务处理。如果客户端永久停止，又没有恢复者接手，这次购买就停在零效果。后续网络和服务恢复可用时，某个参与者需要再次发起处理，才有机会完成购买。

客户端在观察 $L$ 后选择重新 POST，那么在 H1 中也会作出同样选择。H1 已经提交 P1，新 POST 再创建 P2，就多出一次扣款。

固定多等一段时间再重试，也允许构造同样的历史：H1 的响应可以延迟得更久，或永久丢失，而服务随后恢复。一直等待原连接的响应，则需要原连接上仍有可能到达的结果；原请求未执行，或原进程已经退出时，这份结果可能一直缺席。

这个论证约束的是\textbf{没有稳定操作关系、没有新增证据的盲目重发策略}。增加操作身份、持久化结果与恢复入口，可以让后续尝试找到原来的购买。

已有系统采用了这样的设计方向。2015 年提出的 RIFL，英文全称为 Reusable Infrastructure for Linearizability，是一套供服务复用的重复调用处理机制。它给 RPC 分配身份，保存已完成调用的结果，让重试能够找到原记录。Raft 是一种复制共识算法：多台机器各自保存服务数据的一份副本，通过协议确定同一组已提交命令及其执行顺序。Raft 的客户端交互方案另外用调用身份和结果记录处理重复命令\cite{rifl,raft}。这两套设计都为不同请求尝试保留了汇合到原操作的依据。

这段推演只分析当前购买入口中的信息缺口，结论也只覆盖上述重发策略。分布式理论中常提到的 FLP 不可能性结果，研究的则是多进程何时能完成共识\footnote{FLP 取自 Fischer、Lynch、Paterson 三位作者的姓氏。其 1985 年论文研究确定性的二值共识协议：所有已作出的决定必须一致，并且 0 与 1 各自存在可达的决定配置。这项非平凡性要求排除了对所有输入一律决定 0 的程序。在消息可靠但传递时间和进程速度没有已知上界、又允许至多一个进程停止的模型中，满足上述部分正确性条件的协议无法保证每一条允许的执行历史最终都作出决定\cite{flp}。本节采用原文的非平凡性表述；关于购买重发的论证单独来自 H0、H1 两条历史。}。这里两个问题的参与者、目标与故障前提都应分别核对。

至少一次与最多一次还需要指定统计对象。TCP 可能重传数据包，HTTP 入口可能多次受理，数据库可能开始过多笔事务。Alice 关心的是自己授权的一次购买最后产生几笔已提交扣款。恢复过程中可以发生多次尝试，业务规则则要约束它们最后产生的效果。

\subsection*{服务进程退出以后，还剩哪些事实}

先只让购买服务进程崩溃，数据库继续正常运行。服务尚未发送 COMMIT，并且没有其他路径替它提交时，数据库一旦检测到连接终止，就会回滚仍然打开的事务。未提交写入随着事务结束撤回。原连接暂时还没被识别为终止时，事务可能继续占有锁，其他处理也可能等待。

若服务在收到提交成功确认以后退出，恢复进程可以读到 P1、O1、D1。旧进程的响应体已经从内存消失，业务事实仍能由数据库记录重建。恢复者需要知道查询哪个购买编号，查询也需要读到保留这次提交的数据库。

两者之间还有一段特别容易误判：服务已经发出 COMMIT，随后数据库连接失去联系。COMMIT 可能尚未到达，也可能已经处理、确认却没回来。重启使服务能够重新工作，原事务的结果仍要另行查明。

新连接上的 \code{ROLLBACK} 作用于新连接自己的事务。已在旧连接上提交的 P1，仍保留在数据库中。把故障连接移出连接池，再建立新连接，同样只是恢复了通信能力\footnote{连接池保存一组可重复使用的数据库连接，服务处理请求时从中取用、结束后归还。替换故障连接能恢复后续通信，旧连接上那次事务的结果仍须单独查明。}。若捕获提交异常以后立即重新执行整段购买处理，就可能把一次未获知的提交变成第二次购买。

\begin{figure}[htbp]
\centering
\begin{tikzpicture}[x=1cm,y=1cm]
\node[booknode,text width=3.2cm] (a) at (0,0) {尚未发送 COMMIT\\服务退出};
\node[bookresult,text width=3.2cm] (b) at (4.9,0) {已发送 COMMIT\\确认尚未取得时失联};
\node[booknode,text width=3.2cm] (c) at (9.8,0) {已收到成功确认\\服务退出};
\node[booknode,text width=3.2cm] (d) at (0,-2) {数据库检测断连\\回滚仍打开的事务};
\node[bookresult,text width=3.2cm] (e) at (4.9,-2) {命令可能未处理\\也可能已完成提交};
\node[booknode,text width=3.2cm] (f) at (9.8,-2) {数据库保留\\P1、O1、D1};
\node[booknode,text width=3.2cm] (g) at (0,-4) {确认清理结束后\\按未生效路径恢复};
\node[bookresult,text width=3.2cm] (h) at (4.9,-4) {借助操作身份\\查询或继续恢复};
\node[booknode,text width=3.2cm] (i) at (9.8,-4) {读回持久结果\\重新提供给客户端};
\foreach \u/\v in {a/d,b/e,c/f,d/g,e/h,f/i} {\draw[bookarrow] (\u)--(\v);}
\node[booknote,text width=12.4cm] at (4.9,-5.3) {前提：数据库正常运行并保留已确认提交；新连接的 ROLLBACK 作用于新连接自身};
\end{tikzpicture}
\caption{服务退出的位置不同，恢复时可用的证据也不同。提交附近失联需要继续查明原事务的结果。}
\label{fig:ch02-recovery}
\end{figure}


PostgreSQL 18 手册的连接终止说明指出：连接正常或异常终止时，仍然打开的事务会回滚；后台进程也可能先完成正在处理的命令，再发现连接断开\cite{pg}。因此需要结合 COMMIT 的处理位置理解连接失败：事务仍打开时由清理过程回滚；提交已经完成时，数据库保留提交结果。

\subsection*{数据库连接与数据库进程的失败}

连接中断表示双方无法继续沿这条通道通信。数据库进程可以完全健康，P1 也可以完整存在。服务抛出的连接错误描述了通信结果，事务结果则要结合提交或查询证据判断。

数据库进程崩溃，需要数据库恢复机制介入。PostgreSQL 使用 WAL，即 Write-Ahead Logging，通常译为预写日志：先把恢复所需的修改信息写入日志，再按相应规则写入数据文件；崩溃后借助日志恢复已提交修改。

本章对已确认提交的持久性，采用同步提交、可靠存储和正常崩溃恢复等前提。同步提交在本例中指成功确认等待本地事务日志进入持久存储；异步提交则允许提前确认。PostgreSQL 18 手册说明，通常的同步提交会在报告成功前，等待事务 WAL 记录刷入持久存储；异步提交可以先报告成功，最近的事务在崩溃后可能丢失。\code{fsync}、\code{synchronous\_commit} 等设置，以及存储设备兑现写入的情况，都会影响这项保证\cite{pg}。

异步提交、存储失效或恢复旧备份，需要另外分析确认过的记录还保留了哪些。数据库恢复状态内部自洽，与所有已报告成功的购买都仍然存在，也分别对应原子性与持久性的不同要求。

服务日志与数据库日志也各有用途。服务的提交异常说明它没取得成功确认；数据库 WAL 用于数据库自身恢复。购买客户端通常通过授权的业务查询取得结果，直接解释 WAL 则需要数据库层的知识与权限。后续要建立的恢复入口，应当让调用方能定位原购买，并读到与业务效果一致的结果。

\subsection*{客户端崩溃后，谁还记得这次购买}

客户端可能已经把成功响应读入内存，但还没保存 P1，设备就断电了。数据库中已有一次成功购买；重启后的客户端可能只记得用户曾点击购买，也可能连这个动作都没有保留。

客户端已经持久保存 P1，就可以在恢复后查询 P1。若只剩商品编号，还需要区分这次操作、同一账户另一台设备的购买，以及此前迟到的请求。按商品与时间挑一条最近记录，可以帮助人工调查；自动恢复则需要更明确的操作关系。

已显示过成功的客户端，重启后再次显示失败，会让同一次购买产生互相矛盾的说明。保存查询依据、从服务取回结果和继续呈现，需要在客户端恢复流程中接起来。

前面介绍的 RIFL 还把服务器崩溃和数据迁移纳入重复调用处理：执行者重启，或者负责这项数据的服务器改变以后，旧调用的身份和结果仍需要被找回\cite{rifl}。这些用于识别、关联和恢复调用的辅助信息，论文统称为元数据；它们与账户余额等业务数据一起参与系统的恢复。

RIFL 的恰好执行一次保证有明确前提：其底层机制成立、客户端未崩溃，且服务器能可靠保存元数据；论文另行讨论客户端崩溃\cite{rifl}。这项前提与眼前的恢复直接相关：服务器保留完成记录以后，还需要调用方认出自己原先发起的操作。客户端把身份丢失了，恢复流程就必须说明如何重新建立这种联系。

\subsection*{查到一笔购买可以确认什么}

假设维护者从可信关联记录取得 P1，或者客户端已经保存 P1。现在用第一章的详情入口，以 Alice 的身份查询。查询读取主数据库，也就是本例直接接收购买写入的那份数据库，并且假定记录正常保留。

查询采用 PostgreSQL 的 Read Committed 隔离级别，通常译为读已提交。隔离级别规定一个事务能看见其他事务的哪些修改。这里每条普通查询语句都按语句开始时已经提交的数据读取；这份用于判定可见数据的视图称为快照。后续提交可能出现在下一条查询中，本次查询仍按自己的快照读取。

若读到 P1 为 \code{succeeded}，并且关联的 O1、D1 完整、金额为 3000，就确认了这笔购买已经提交。购买、订单与扣款在本章规则下由同一事务保存，查询通过新的通信路径把这些事实传给调用方。

这份结果说明 P1 的情况。如果 R2 已经创建了 P2，P2 可能同时存在；核对一次购买意图是否多扣，需要把相关尝试与全部对应流水联系起来。查到某一笔成功，与检查相关操作只有这一笔，分别需要存在证据与完整范围内的核对。

查到 \code{insufficient\_funds}，说明 P1 在当时得到余额不足的终态，订单和扣款为空。后来充值改变当前余额，P1 仍保留原来的处理结果。用户此后决定购买，可以发起新操作。

详情入口按账户与购买编号取记录：以 Alice 的身份查询 P1，就读取属于她的 P1；查询 P2，取得的是 P2；其他账户的编号会受到授权过滤。200、404 等 HTTP 状态码，需要连同查询使用的账户、编号和返回的购买状态理解。

\subsection*{没查到为什么经常还不能重发}

在另一条历史中，恢复者已知 P1，原事务 T1 却仍未提交。T1 已在内部写入 P1 并扣款，另一个连接执行普通 \code{SELECT}，此时只读到已提交数据，返回无记录。T1 稍后提交，P1 才成为可见的成功购买。

PostgreSQL 的 Read Committed 普通查询读取语句开始时的快照。如果 P1 在查询开始以后才提交，本次查询也可能读不到它\cite{pg}。有些系统会把主数据库的修改随后传给另一份数据库供查询使用，这种异步副本在收到并应用修改以前，会暂时缺少主库已有的新记录；这段差距叫复制滞后。本例直接读主库，消除了这项差距，进行中的事务与查询之间仍存在这种交错。

原请求还可能在开始事务以前，就排在服务队列里。维护者查询最近购买，结果为空，于是再次 POST。第二次请求先提交 P2，队列中的旧请求稍后执行，又提交 P1。那次空查询准确反映了当时的数据，旧请求以后仍有执行机会。

\begin{figure}[htbp]
\centering
\begin{tikzpicture}[x=1cm,y=1cm]
\foreach \x in {0,4,8,12} {\draw[Ink!35,dashed] (\x,-.5)--(\x,-6.2);}
\node[booknode,text width=1.6cm] at (0,0) {旧请求 R1};
\node[booknode,text width=1.6cm] at (4,0) {恢复查询};
\node[booknode,text width=1.6cm] at (8,0) {重试 R2};
\node[booknode,text width=1.6cm] at (12,0) {主数据库};
\node[bookresult,text width=2cm] at (0,-1.1) {执行前暂停};
\draw[bookarrow] (4,-1.8)--node[above,font=\scriptsize]{查询购买}(12,-1.8);
\draw[bookarrow] (12,-2.5)--node[above,font=\scriptsize]{无记录，余额 10000}(4,-2.5);
\draw[bookalert] (8,-3.3)--node[above,font=\scriptsize]{P2 提交}(12,-3.3);
\node[bookresult,text width=2cm] at (12,-4.1) {余额 7000};
\draw[bookalert] (0,-4.9)--node[above,font=\scriptsize]{R1 恢复执行，P1 也提交}(12,-4.9);
\node[bookresult,text width=2.6cm] at (12,-5.8) {余额 4000\\两条扣款流水};
\node[booknote,text width=7cm] at (4.5,-6.7) {\mbox{空查询准确反映当时的数据}\\\mbox{旧请求仍有后续执行机会}};
\end{tikzpicture}
\caption{查询无记录以后，重试与旧请求先后提交。普通只读查询没有改变原处理的执行资格。}
\label{fig:ch02-query-race}
\end{figure}


把旧请求暂停在购买逻辑执行前，客户端超时后，主库新连接仍会读到余额 10000、购买记录为空。允许重试先完成，余额变为 7000；再释放旧请求，原处理也可以完成，余额变为 4000，留下两条流水。延迟位于业务执行以前，或位于事务已经写入但尚未提交时，都会给空查询以后的旧执行留下机会。

网络中尚未送达的请求、服务队列里的请求、尚未提交的事务，都可能在空查询以后继续生效。要在无记录以后重新开始处理，需要进一步确认原执行已结束并回滚，而且没有其他原尝试仍可生效；或者通过协议把后续尝试与原尝试汇合到同一操作。

异步副本、旧快照、错误编号和授权过滤，会进一步影响空查询的含义。第一章的 404 同时用于编号不存在和无权访问，调用方要按详情入口的受理规则理解它。

章首的客户端还没有取得 P1，因而连按 P1 查询这一步都缺少依据。增加按商品编号查询最近订单的入口，可以帮助展示历史；恢复原购买仍然需要识别它与用户这次动作的关系。

\subsection*{重试、查询、等待和取消各自承担什么代价}

请求未到达或事务已经回滚后，重试可以推动购买继续完成。在当前入口中，它也可能创建第二笔购买，并增加服务负载。重试退避是让连续重试之间留出等待间隔，例如失败次数增加时逐步延长间隔；它和次数限制一起控制请求压力。识别同一操作则需要另一项机制。

查询读取已有事实。调用方拥有正确购买身份，结果也能读取时，可以借此缩小未知范围。查到确定终态以后，就能据此结束恢复；尚无记录时，还需处理原执行继续生效的可能。客户端隔一段时间主动再查一次，称为轮询。查询频率影响成本，也影响客户端多快得知结果。

等待给延迟响应和进行中的事务留出时间。服务最终完成、结果仍可取得时，等待可能省去一次重新处理。原进程已退出、旧响应永久丢失时，等待原连接就缺少进展条件。界面可以在等待上限后进入待确认状态，再由恢复流程继续查询和处理。

取消要影响远端执行，需要取消消息到达执行者，并且与业务提交有明确的处理顺序。关闭页面、停止等待和断开 HTTP 连接，发生在客户端这一侧；服务的事务可能仍在执行，也可能已经提交。

\begin{figure}[htbp]
\centering
\begin{tikzpicture}[x=1cm,y=1cm]
\node[font=\small\sffamily\bfseries] at (2,1.1) {取消先被执行者处理};
\node[font=\small\sffamily\bfseries] at (9,1.1) {购买提交先完成};
\node[booknode,text width=4.7cm] (a) at (2,0) {事务尚未提交};
\node[booknode,text width=4.7cm] (b) at (2,-1.5) {执行者取消并回滚\\调用方取得回滚结果};
\node[booknode,text width=4.7cm] (c) at (2,-3) {购买无效果\\余额 10000};
\draw[bookarrow] (a)--(b);
\draw[bookarrow] (b)--(c);
\node[booknode,text width=4.7cm] (d) at (9,0) {P1、O1、D1 提交\\余额 7000};
\node[bookresult,text width=4.7cm] (e) at (9,-1.5) {取消消息随后到达\\原购买已完成};
\node[booknode,text width=4.7cm] (f) at (9,-3) {查询原结果\\返还另建退款操作};
\draw[bookalert] (d)--(e);
\draw[bookarrow] (e)--(f);
\node[booknote,text width=11.7cm] at (5.5,-4.2) {客户端停止等待发生在本地；远端取消需要执行者处理，并取得明确业务结果};
\end{tikzpicture}
\caption{取消与提交之间需要仲裁顺序。已提交购买保留原记录，后续返还另记一项操作。}
\label{fig:ch02-cancel-race}
\end{figure}


PostgreSQL 18 手册的取消请求说明提供了一个具体对照：取消进行中的查询使用另一条连接，取消请求没有直接的成功答复；信号到达时，原查询也可能已经完成，客户端还要等待原命令结果\cite{pg}。

即使数据库取消确实影响了当前事务，业务层仍要确认回滚和其他尝试的情况。购买已经提交时，返还款项需要建立退款或撤销操作，并保留原扣款历史。这种针对已完成业务另外发起修正动作的做法称为补偿；退款是新的业务动作，原扣款仍然曾经发生。

\begin{center}
\begin{minipage}{\linewidth}
\centering\small\captionsetup{hypcap=false}
\captionof{table}{超时以后几种动作的作用}
\begin{tabularx}{\linewidth}{p{.19\linewidth}XX}
\toprule
\textbf{动作} & \textbf{推进所需的条件} & \textbf{要继续处理的问题} \\
\midrule
再次提交 & 原请求未生效，服务恢复可用 & 原请求已完成或仍在途，可能再次生效 \\
查询购买 & 已知身份，结果记录可读取 & 无记录以后旧执行仍可能提交 \\
等待或延后查询 & 原执行能够结束，结果仍可取得 & 永久等待、资源占用与再次点击 \\
请求取消 & 远端能把取消和执行有序处理 & 取消晚到，已提交购买需要补偿 \\
\bottomrule
\end{tabularx}
\end{minipage}
\end{center}

界面中的取消按钮如果只终止网络等待，可以显示停止等待；宣告购买已取消，则需要远端业务结果支撑。每个动作作用于哪笔操作、可以改变哪种状态、失败后留下什么证据，都需要由流程明确。

最多一次扣款要求所有允许历史都满足效果上限；最终取得终态则需要网络、服务、数据库与恢复者在某些时段能够继续工作。客户端永久停止，可以减少再次扣款的机会，却会让 H0 中的购买一直没有发生。持续创建新购买，很容易取得某次成功，也可能多次扣款。后续设计需要同时说明效果约束与继续推进的条件。
\FloatBarrier
